% Author: Charis Opol (F-11). Draft of 23 Nov 2026. Every finding below is taken from the paper's
% abstract or full text in "final report/" (PDFs), except Zhou 2019, Zhang 2019 and Chung 2017,
% which are cited from their titles and the reading list only (see the author queries).
\section{Related Work}
\label{sec:related}

\subsection{Lip reading}
Neural lip reading began on GRID. Wand et al.~\cite{wand2016} stacked feed-forward and LSTM
layers into one network trained end to end and reached 79.6~\% word accuracy on held-out GRID
speakers, 11.6 points above the best system built on hand-crafted features. LipNet~\cite{assael2016}
moved from words to whole sentences: spatiotemporal convolutions, a recurrent network and a CTC
loss~\cite{graves2006} gave 95.2~\% sentence-level accuracy on GRID, although with speakers shared
between training and test. LCANet~\cite{xu2018} added a cascaded attention-CTC decoder and reported
3.0~\% WER on GRID. Outside GRID, Stafylakis and Tzimiropoulos~\cite{stafylakis2017} combined a
3-D convolution, a ResNet applied to each frame and a bidirectional LSTM, and reached 83.0~\% word
accuracy on the 500-word Lip Reading in the Wild benchmark. Martinez et al.~\cite{martinez2020}
replaced the recurrent layers of that design with temporal convolutional networks and simplified
its training. LipScribe's lip encoder follows the front end of~\cite{stafylakis2017}, and our
lips-only GRID baselines are tested, as in~\cite{wand2016}, on speakers not seen in training.

\subsection{End-to-end audio-visual speech recognition}
Petridis et al.~\cite{petridis2018e2e} trained the first audio-visual model that learns from raw
pixels and raw waveforms together, with one ResNet and recurrent stream per modality. On clean
audio it was only slightly better than an audio-only model, but at high noise levels it was much
better. The same pattern held at sentence level. A hybrid CTC/attention model on LRS2 was 1.3 points
better than audio alone on clean speech and up to 32.9 points better in noise~\cite{petridis2018hybrid},
and Afouras et al.~\cite{afouras2022} found that their audio-visual transformers outperformed both
single-stream models under loud babble noise, with the CTC variants coping better with noise than
the sequence-to-sequence ones. Conformer encoders over ResNet-18 features then raised the state of
the art on LRS2 and LRS3~\cite{ma2021}, and Auto-AVSR~\cite{ma2023} reached 0.9~\% WER on LRS3 by
adding automatically transcribed VoxCeleb2~\cite{chung2018voxceleb2} and AVSpeech clips to the
training data. Self-supervised pre-training lowers the need for labels: with AV-HuBERT, Shi
et al.~\cite{shi2022robust} halved the babble-noise WER of the previous state of the art on LRS3
(28.0~\% to 14.1~\%) using 30 hours of labelled data rather than 433. Whisper-Flamingo~\cite{rouditchenko2024}
injects lip features into the Whisper recogniser~\cite{radford2023whisper} through gated
cross-attention and improves on audio-only Whisper in noise. Our real-world model builds on two of
these lines: its lip encoder is that of Auto-AVSR, and its audio encoder is the self-supervised
WavLM~\cite{chen2022wavlm}.

\subsection{Fusion, gating and modality reliability}
How to combine the two streams is the central design choice, and three findings shaped ours. First,
a fused model can be \emph{worse} than audio alone on clean speech. Tao and Busso~\cite{tao2018}
observed this drop in a hybrid DNN-HMM system and added a gating layer that suppresses uninformative
visual features, which kept or improved accuracy in every condition they tested. Second, the right
weight for each stream changes with the noise. Meutzner et al.~\cite{meutzner2017} derived dynamic
stream weights from an estimate of audio reliability and outperformed a fixed weighting that knew
the true SNR, and Zhou et al.~\cite{zhou2019} proposed modality attention for end-to-end AVSR. More
recently, AV-RelScore~\cite{hong2023} scores the reliability of each stream so that a model can rely
on the cleaner one even when the video is also corrupted, and router-gated cross-attention~\cite{lim2025}
down-weights audio tokens that a router judges corrupted, cutting WER by 16.5 to 42.7~\% relative to
AV-HuBERT on LRS3. Third, a model trained on both streams can learn to ignore the weaker one.
Sterpu et al.\ showed that attention-based fusion learns to align the two streams and improves on
audio alone by 7 to 30~\% on TCD-TIMIT~\cite{sterpu2018}, but that on the harder LRS2 it at first
gave no improvement at all, which they addressed with an auxiliary lip task~\cite{sterpu2020}.
Multi-condition training with modality dropout~\cite{zhang2019} is a direct way to stop a network
leaning on one stream. Abdelaziz~\cite{abdelaziz2018} compared five classical fusion models on one
benchmark and is a reference point for the concatenation baselines we keep. Model P combines these
ideas: cross-modal attention in the spirit of~\cite{sterpu2018}, a per-frame gate as
in~\cite{tao2018}, and modality dropout as in~\cite{zhang2019}.

\subsection{Measuring what the lips contribute}
A low WER does not show that a model uses the lips. Lin and Harte~\cite{lin2025} measured the
\emph{effective SNR gain}, the drop in SNR at which the audio-visual model matches the audio-only
model's WER at 0~dB, for three recent systems. Auto-AVSR had the best WER on every noise level but
the smallest gain, 3.7~dB on LRS2 and 2.5~dB on LRS3, and the authors concluded that current systems
do not fully exploit the visual stream. We therefore compare each fused model with an audio-only
model trained in the same way, test what happens when one stream is removed or replaced, and report
the effective SNR gain alongside WER (Sections~\ref{sec:results} and~\ref{sec:realworld}).

\subsection{Noise, data and synchrony}
NTCD-TIMIT~\cite{abdelaziz2017} established the practice we follow on GRID: add several noise types
to clean audio-visual speech at a range of SNRs, and give baselines for audio-only, video-only and
audio-visual models. Our noise comes from MUSAN~\cite{snyder2015}, a freely licensed corpus of music,
speech and noise. Our real-world clips come from TED talks, the source of LRS3~\cite{afouras2018},
and from films, and the Auto-AVSR checkpoint we start from was trained on LRS3 and VoxCeleb2. Finally,
SyncNet~\cite{chung2017} learns whether a voice and a face are in sync, which is the kind of signal
our negative result in Section~\ref{sec:realworld} calls for. Ivanko et al.~\cite{ivanko2023} survey
the deep-learning AVSR work of 2013 to 2023 in more breadth than we can here.
